\section*{الفصل \ref{ch:g2:moon-in-the-sky} --- القمر في السماء}

\begin{solution}{exo:g2:moon-in-the-sky:1}
أصغر من الأرض بكثير، وأبعد من الغيوم بكثير --- أبعد من أن تبلغه
طائرة. ونعلم أن الناس يستطيعون زيارته لأن روّاد الفضاء فعلوا:
واستغرقت الرحلة ثلاثة أيام في كل اتجاه.
\end{solution}

\begin{solution}{exo:g2:moon-in-the-sky:2}
الفوّهات --- ندوب على شكل أوعية صنعتها صخور من الفضاء اصطدمت
\omterm{def:g2:moon-in-the-sky:moon}{بالقمر}. وهي باقية لأن \omterm{def:g2:moon-in-the-sky:moon}{القمر} ليس عنده \omterm{def:g2:air-around-us:wind}{ريح} ولا مطر يبليانها.
\end{solution}

\begin{solution}{exo:g2:moon-in-the-sky:3}
لاقط \omterm{def:g1:light-and-shadows:source}{ضوء}. \omterm{def:g1:light-and-shadows:source}{فضوء} \omterm{def:g2:moon-in-the-sky:moon}{القمر} \omterm{def:g1:light-and-shadows:source}{ضوء} شمس يرتدّ عن صخر \omterm{def:g2:moon-in-the-sky:moon}{القمر} الرمادي، كما
يضيء الجدار الأبيض غرفةً من غير أن يكون مصباحًا.
\end{solution}

\begin{solution}{exo:g2:moon-in-the-sky:4}
النصف المواجه للشمس دائمًا --- فالجهة المشمسة ساطعة، والجهة
البعيدة في ليلها الخاص.
\end{solution}

\begin{solution}{exo:g2:moon-in-the-sky:5}
نعم --- ففي كثير من العصريات يتعلّق \omterm{def:g2:moon-in-the-sky:moon}{قمر} شاحب في السماء الزرقاء.
وخير برهان لصديق شاكّ: أن تخرجا في عصريّة كهذه وتشير إليه.
\end{solution}

\begin{solution}{exo:g2:moon-in-the-sky:6}
\omterm{ex:g2:moon-in-the-sky:shapes}{هلال}، فنصف، فيكاد يكتمل، فبدر --- ثم هبوط راجع عبر ما يكاد
يكتمل، فالنصف، \omterm{ex:g2:moon-in-the-sky:shapes}{فالهلال}، إلى المحاق المختبئ. ويستغرق العرض كله
نحو شهر.
\end{solution}

\begin{solution}{exo:g2:moon-in-the-sky:7}
لا تُحدث الصرخة صوتًا --- فالصوت يحتاج إلى \omterm{def:g2:air-around-us:air}{هواء}. ولا تستطيع
الطائرة الورقية أن تطير: فانعدام \omterm{def:g2:air-around-us:air}{الهواء} يعني انعدام \omterm{def:g2:air-around-us:wind}{الريح}، ولا
شيء تستند إليه الطائرة.
\end{solution}

\begin{solution}{exo:g2:moon-in-the-sky:8}
الساعة نفسها، والرسم الأمين نفسه في كل مساء (والمكان نفسه يساعد
أيضًا): فالمفكّرة العادلة لا تغيّر إلا شيئًا واحدًا في كل مرة
--- الليلة.
\end{solution}

\begin{solution}{exo:g2:moon-in-the-sky:9}
البيوت قريبة، فمرورها في السيارة يغيّر منظرها عند روزا ---
تنزلق وتذهب. أما \omterm{def:g2:moon-in-the-sky:moon}{القمر} فمن البعد بحيث لا تغيّر الرحلة كلها
منظره لها البتة: فيبقى عند كتفها كأنه تابع. وبطل البعد ``يتبع''
دائمًا.
\end{solution}

\begin{solution}{exo:g2:moon-in-the-sky:10}
الأرض نفسها. \omterm{def:g1:light-and-shadows:source}{فضوء} الشمس يرتدّ عن الأرض الساطعة إلى جانب \omterm{def:g2:moon-in-the-sky:moon}{القمر}
المعتم فيضيئه إضاءة خافتة --- وهو \omterm{def:g1:light-and-shadows:source}{الضوء} الأرضي: كوكبنا يردّ
\omterm{def:g2:moon-in-the-sky:moon}{للقمر} جميله.
\end{solution}
